% !PRINTED none
\begin{center}
\gesperrt{Third Chapter.}

\gesperrt{The Group of the Cyclotomic Fields.}

\rule{3cm}{0.4pt}

\S.\,17.

\gesperrt{The Resolvents of the Theory of the Division of the Circle.}
\end{center}

Of the theorems on \gesperrt{Abelian} groups we make an
application to the theory of the division of the circle for the case in which the
degree of the roots of unity is not a prime, but a
higher power of a prime. If $q$ is an odd prime
and $m = q^{\varkappa}$, $\varkappa > 1$, we take a primitive congruence root
 $g$ of $m$ and set, for every number $n$ not divisible
 by $q$,
\weq{1}{n \equiv g^{\nu} \pmod{m}.}
Then $\nu$ is the \gesperrt{index} of $n$. If $n$ runs through the group $\mathfrak{N}$
of the number classes not divisible by $q$ with respect to the modulus $m$,
of degree
\weq{2}{c = \varphi(q^{\varkappa}) = q^{\varkappa - 1}(q - 1),}
then $\nu$ runs through a complete system of residues with respect to the modulus $c$. We
now take a primitive $m$th root of unity $r$ and a primitive
 $c$th root of unity $\varepsilon$ and form the \gesperrt{Lagrange}
\gesperrt{resolvents}
\weq{3}{(\varepsilon^{\beta}, r) = \sum^{n} \varepsilon^{\beta \nu} r^n,}
in which $\beta$ is an arbitrary exponent. The question at issue is
when such a resolvent can vanish. If
$\varkappa = 1$, so that $m$ is a prime, it vanishes, as we saw in
Bd.\ I, \S.\,169, for no value of $\beta$. In the general
 case of an arbitrary $\varkappa$, let $n'$ denote an arbitrary number
from $\mathfrak{N}$. Then $n n'$ runs through the whole group $\mathfrak{N}$ together with $n$.
One can therefore replace $n$ by $n n'$ in (3), if at the same time
%CONT>>
